\subsection{Реализация базы данных}
\label{sec:development:database}

\subsubsection{Физическая схема}

Физическая схема получена из логической (рисунок~\ref{fig:design:logical}) по обычным правилам: каждая сущность стала таблицей, каждый атрибут~-- столбцом, связи «один ко многим» выражены внешними ключами в таблице на стороне «многие». Физическая схема базы данных приведена на рисунке~\ref{fig:development:physical}, назначение таблиц и характер их роста~-- в таблице~\ref{table:development:tables}.

\begin{figure}[H]
	\centering
	\includegraphics[width=\linewidth]{kufar/physical}
	\caption{Физическая схема базы данных}
	\label{fig:development:physical}
\end{figure}

\begin{table}[ht]
\caption{Таблицы базы данных}
\label{table:development:tables}
\centering
	\begin{tabular}{|>{\raggedright}m{0.27\textwidth}
	                |>{\raggedright}m{0.38\textwidth}
	                |>{\raggedright\arraybackslash}m{0.25\textwidth}|}
	\hline
	\centering Таблица & \centering Назначение & \centering\arraybackslash Рост \\
	\hline Пользователи & Учётные записи, настройки, электронная почта & По одной строке на пользователя \\
	\hline Отслеживаемые товары & Критерии поиска пользователей & Не более 25 строк \\
	\hline Объявления & Сведения, полученные с площадки & С каждым циклом опроса, очищается через 30 дней \\
	\hline Совпадения & Подошедшие объявления и очередь доставки & По мере отбора \\
	\hline Журнал уведомлений & Факты отправки с ценой & С каждой отправкой, очищается через 90 дней \\
	\hline Атрибуты фильтров & Характеристики категорий площадки & Пересобирается раз в неделю \\
	\hline Значения атрибутов & Возможные значения характеристик & Пересобирается раз в неделю \\
	\hline Пользовательские атрибуты & Резерв для расширения системы & Не заполняется \\
	\hline
	\end{tabular}
\end{table}

Типы данных выбраны по смыслу атрибутов. Первичными ключами всех таблиц служат целые числа с автоматическим приращением. Идентификатор пользователя \textit{Telegram} хранится 64-битным целым: значения таких идентификаторов уже выходят за пределы 32-битного типа. Строки с известной верхней границей длины, например имя, код характеристики или адрес электронной почты, хранятся в строковых полях ограниченной длины, до 255 символов. Заголовок, описание и ссылка объявления не имеют разумного предела и хранятся в текстовых полях без ограничения длины. Хеш кода подтверждения занимает ровно 64 символа, по длине шестнадцатеричной записи результата \textit{SHA-256}.

Цены хранятся числами с плавающей точкой двойной точности. В системе цены только сравниваются с порогом и между собой, а доля выгоды округляется до четырёх знаков, поэтому погрешность представления на результат не влияет. Если в систему будут добавлены суммирование или денежные расчёты, цены следует перевести в десятичный тип с фиксированной точностью.

Моменты времени хранятся без часового пояса и всегда в \textit{UTC}. Время публикации, которое площадка передаёт с указанием пояса, система приводит к \textit{UTC} перед записью. Благодаря этому сравнение «объявление опубликовано после добавления товара» корректно при любой настройке часового пояса сервера. Для дат создания и изменения записей значение по умолчанию задаёт сама СУБД текущим временем.

Поля с фиксированным набором значений (категория товара, площадка, реакция пользователя) объявлены перечислимыми типами \textit{PostgreSQL}. СУБД отклоняет любое значение вне перечня, и опечатка в коде не может записать в базу несуществующую категорию. Признаки хранятся логическим типом, а счётчики попыток целыми числами.

Тип \textit{JSON} использован в семи столбцах: обязательные слова, слова-исключения и фильтры площадки у товара, фотографии, продавец и исходный ответ у объявления, детали события у совпадения. У этих данных общее свойство: их всегда читают и записывают целиком и никогда не ищут в базе по отдельному элементу. Фильтры площадки, например, подставляются в запрос к \textit{Kufar.by} в неизменном виде.

\subsubsection{Нормализация}

Схема приведена к третьей нормальной форме с несколькими осознанными отступлениями.

Первая нормальная форма требует атомарности значений. Формально её нарушают \textit{JSON}-столбцы: список слов-исключений или набор фильтров хранится в одной ячейке. Вынесение их в отдельные таблицы дало бы одну-две лишние таблицы на каждый список и соединение при каждом чтении товара, но ни одного нового запроса, которому такая структура была бы нужна. Система не ищет товары по слову-исключению и не считает, сколько пользователей выбрали бренд. Когда такие запросы появятся, фильтры можно будет перенести в отдельную таблицу связи между товаром и значением атрибута без изменения остальной схемы.

Вторая нормальная форма выполняется автоматически: все таблицы имеют простой суррогатный первичный ключ, и частичной зависимости от части ключа возникнуть не может.

Третья нормальная форма требует, чтобы неключевые атрибуты не зависели друг от друга. Это требование намеренно нарушено в журнале уведомлений. Площадка, внешний идентификатор объявления и цена в нём дублируют данные таблицы объявлений. Дублирование сделано по двум причинам. Проверка повтора (отправлялось ли это объявление пользователю и по какой цене) выполняется для каждого уведомления, и без копии каждая такая проверка требовала бы соединения с таблицей объявлений. Кроме того, в журнале записана цена на момент отправки, а в таблице объявлений хранится текущая цена, которая со временем меняется. Это разные факты, поэтому формально зависимость отсутствует, хотя названия столбцов совпадают.

Данные подтверждения электронной почты (хеш кода, срок, число попыток) хранятся в таблице пользователей, хотя по смыслу образуют отдельный объект. Пользователь связан с ними отношением один к одному, данные невелики и всегда читаются вместе с пользователем, поэтому отдельная таблица только добавила бы соединение.

\subsubsection{Индексы}

Каждый индекс ускоряет чтение и замедляет запись: при вставке и обновлении строки СУБД перестраивает все индексы таблицы. Таблица объявлений обновляется на каждом цикле опроса, поэтому индексы подбирались не по столбцам, а по запросам, которые система действительно выполняет. Все индексы~-- сбалансированные деревья (\textit{B-tree}), принятые в \textit{PostgreSQL} по умолчанию: они обслуживают и поиск по равенству, и сравнение по диапазону, которое нужно при подсчёте уведомлений за период. Соответствие индексов запросам приведено в таблице~\ref{table:development:indexes}.

\begin{table}[ht]
\caption{Индексы и обслуживаемые ими запросы}
\label{table:development:indexes}
\centering
	\begin{tabular}{|>{\raggedright}m{0.3\textwidth}
	                |>{\raggedright}m{0.37\textwidth}
	                |>{\raggedright\arraybackslash}m{0.23\textwidth}|}
	\hline
	\centering Индекс & \centering Запрос & \centering\arraybackslash Частота \\
	\hline Объявления: площадка и внешний идентификатор (составной) & Поиск сохранённого объявления при обработке ответа площадки & До 1\,050 раз за цикл \\
	\hline Пользователи: идентификатор \textit{Telegram} (уникальный) & Определение пользователя по входящему сообщению & Каждое сообщение \\
	\hline Товары: владелец & Список товаров пользователя, подсчёт занятости личного лимита & Каждый диалог и каждое добавление \\
	\hline Совпадения: товар & Выборка очереди доставки, отчёт за период & Каждую минуту \\
	\hline Совпадения: объявление & Проверка, есть ли у объявления совпадения, при очистке & Раз в сутки \\
	\hline Журнал: получатель & Подсчёт уведомлений пользователя за сутки & Каждая отправка \\
	\hline Журнал: товар & Подсчёт уведомлений по товару за час & Каждая отправка \\
	\hline Журнал: внешний идентификатор объявления & Проверка, отправлялось ли объявление и по какой цене & Каждая отправка \\
	\hline Справочник: категория площадки, атрибут & Построение меню фильтров, обновление справочника & Каждое открытие фильтров \\
	\hline
	\end{tabular}
\end{table}

Самый частый запрос~-- поиск объявления по паре «площадка~-- внешний идентификатор». Поиск выполняется для каждого объявления в каждом цикле: при 25 товарах и 42 объявлениях на запрос это до 1\,050 поисков за цикл и более 500\,000 в сутки. Для него создан составной индекс по двум столбцам. До его появления на этих столбцах были только отдельные индексы, и планировщику запросов приходилось объединять их результаты. Порядок столбцов в составном индексе выбран так, чтобы индекс годился и для запросов по одной площадке: ведущим стоит столбец площадки, а внешний идентификатор уточняет запись внутри неё. Поскольку площадка в системе пока одна, отдельный индекс по ней стал избыточным: выборка не сужается, а запись замедляется, поэтому такой индекс можно удалить.

Уникальный индекс по идентификатору \textit{Telegram} решает две задачи сразу. Индекс поддерживает ограничение уникальности и обслуживает запрос, который выполняется чаще всех остальных запросов бота: на каждое входящее сообщение система находит пользователя по идентификатору его учётной записи. Без индекса этот поиск был бы последовательным просмотром всей таблицы пользователей.

Индекс по владельцу в таблице товаров обслуживает два запроса: вывод списка товаров пользователя и подсчёт занятости личного лимита перед добавлением нового товара. Оба выполняются с условием равенства по владельцу, и выборка по индексу возвращает сразу нужные несколько строк.

В таблице совпадений проиндексированы оба внешних ключа, но обслуживают они разные запросы. Индекс по товару соединяет совпадения с товарами при выборке очереди доставки и при формировании отчёта за период. Индекс по объявлению нужен задаче очистки: перед удалением устаревшего объявления система проверяет, не ссылается ли на него хотя бы одно совпадение, и без индекса такая проверка означала бы просмотр всей таблицы совпадений для каждого удаляемого объявления.

Три индекса журнала уведомлений поддерживают антиспам. По получателю считаются уведомления пользователя за текущие сутки, по товару~-- уведомления за последний час, по внешнему идентификатору объявления проверяется, отправлялось ли объявление этому пользователю и по какой цене. Все три запроса выполняются перед каждой отправкой, поэтому их стоимость напрямую влияет на длительность цикла доставки. Запрос проверки повтора отбирает строки по получателю, площадке и внешнему идентификатору одновременно; сейчас планировщик использует один из индексов и проверяет остальные условия по найденным строкам, а при росте журнала выгоднее станет составной индекс по получателю и внешнему идентификатору.

Справочник фильтров индексируется по категории площадки в таблице атрибутов и по атрибуту в таблице значений. Эти индексы обслуживают построение меню фильтров в боте, а уникальные ограничения на пару «категория~-- системное имя» и пару «атрибут~-- код значения» дополнительно ускоряют еженедельное обновление справочника: при повторном сборе система ищет существующую запись именно по этим парам.

Индексов по цене и категории в таблицах объявлений и товаров нет, и это сделано сознательно. Отбор объявлений по цене и категории выполняет сама площадка: система передаёт порог и категорию в поисковом запросе и получает только подходящие объявления. В базе данных по цене объявления не ищут. Таблица товаров содержит не более 25 строк, и последовательный просмотр такой таблицы быстрее обращения к индексу.

При росте нагрузки естественным улучшением станет частичный индекс по совпадениям, ещё не доставленным пользователю. Очередь доставки выбирает только их, а таких записей в любой момент единицы, тогда как в таблице совпадений хранится вся история. Частичный индекс занимает место, пропорциональное числу недоставленных совпадений, и не растёт вместе с историей.

\subsubsection{Ограничения целостности}

Целостность данных обеспечивается ограничениями на уровне СУБД и проверками на входе в серверный интерфейс.

На уровне СУБД действуют:
\begin{itemize}
	\item первичные ключи во всех таблицах;
	\item внешние ключи с явно заданными правилами удаления: товар ссылается на пользователя, совпадение на товар и объявление, запись журнала на пользователя, товар, объявление и совпадение, значение атрибута на атрибут;
	\item ограничения уникальности: идентификатор \textit{Telegram} пользователя, пара «категория~-- системное имя» у атрибута фильтра, пара «атрибут~-- код значения» у значения;
	\item ограничения \textit{NOT NULL} на всех обязательных столбцах; необязательными оставлены только данные, которых действительно может не быть: цена договорного объявления, адрес почты, время доставки недоставленного совпадения;
	\item перечислимые типы для категорий, площадок и реакций;
	\item значения по умолчанию на стороне СУБД для дат создания и изменения записей и для столбцов, добавленных миграциями; начальные значения признаков задаёт приложение при вставке записи.
\end{itemize}

Ограничения уникальности справочника позволяют пересобирать его без дублей: при повторном сборе существующее значение обновляется, а не добавляется второй раз. Совпадение для пары «товар~-- объявление» создаётся только после проверки, что такого совпадения ещё нет. Проверка и вставка выполняются в одной транзакции цикла опроса.

Удаление товара затрагивает зависимые записи, поэтому правило удаления задано в схеме для каждого внешнего ключа. Совпадения, пользовательские атрибуты и записи журнала уведомлений по товару удаляются вместе с ним (\textit{ON DELETE CASCADE}). Ссылка журнала на совпадение обнуляется (\textit{ON DELETE SET NULL}): запись журнала хранит собственную копию цены и внешнего идентификатора объявления и без совпадения не теряет смысла. По тому же правилу вместе с пользователем удаляются его товары и журнал, а вместе с атрибутом фильтра~-- его значения. У ссылки журнала на объявление правила удаления нет: очистка удаляет только объявления без совпадений, а на них журнал не ссылается. Правила выполняет сама СУБД, поэтому одного запроса \textit{DELETE} к таблице товаров достаточно, чтобы удалить товар со всеми зависимыми записями в одной транзакции.

Проверки, которые неудобно выразить ограничениями СУБД, выполняются при приёме данных серверным интерфейсом: пороговая цена больше нуля и не больше 10~000~000 рублей, название не длиннее 255 символов, не более 20 слов в каждом списке и не более 100 символов в слове, коды фильтров принадлежат категории товара. Файл \textit{CSV} при импорте ограничен 256~КБ и 100 строками. Некорректный запрос отклоняется с понятным сообщением до обращения к базе данных.

Каждый запрос к серверному интерфейсу выполняется в одной транзакции: при успехе изменения фиксируются, при любой ошибке транзакция откатывается целиком. Цикл опроса площадки записывает объявления и совпадения всех товаров одной транзакцией. Доставка уведомлений отмечает совпадение доставленным и добавляет запись в журнал в одной транзакции, поэтому отметка без записи или запись без отметки появиться не могут.

\subsubsection{Миграции и обслуживание данных}

Схема базы данных изменялась пятью миграциями. Первая создаёт все таблицы и индексы. Вторая добавляет в таблицу совпадений счётчик попыток доставки и текст ошибки, на которых построена повторная отправка. Третья добавляет составной индекс по площадке и внешнему идентификатору объявления. Четвёртая добавляет поля электронной почты в таблицу пользователей и отдельную почтовую очередь в таблицу совпадений. Пятая задаёт правила удаления для внешних ключей и удаляет избыточный индекс по площадке. Новые обязательные столбцы добавлялись со значением по умолчанию на стороне СУБД, поэтому миграции применялись к таблицам, уже содержащим данные, без их пересоздания.

Миграции применяет серверный сервис при запуске. Перед применением сервис берёт рекомендательную блокировку \textit{PostgreSQL}: если одновременно запускаются два экземпляра сервиса, второй дожидается, пока первый закончит миграцию, и находит схему уже в актуальном состоянии.

Объявления сохраняются на каждом цикле опроса, и без очистки их таблица росла бы всё время работы системы. Раз в сутки удаляются объявления, которые не появлялись в выдаче 30 дней и не участвуют ни в одном совпадении. Объявления с совпадениями сохраняются, потому что на них ссылается история уведомлений и отчёты пользователя. Записи журнала уведомлений старше 90 дней удаляются по тому же расписанию.

\subsubsection{Резервное копирование}

Резервное копирование выполняет отдельный контейнер. Раз в сутки контейнер создаёт полный логический дамп базы данных утилитой \textit{pg\_dump}, сжимает его и сохраняет в каталог на сервере вне тома СУБД. Хранятся семь последних копий, более старые удаляются автоматически.

Резервная копия, которую ни разу не восстанавливали, не даёт гарантии, что данные удастся вернуть. Поэтому в систему входит сценарий проверки восстановления. Сценарий создаёт временную базу данных, загружает в неё последний дамп, подсчитывает строки в основных таблицах и удаляет временную базу. Проверка запускается вручную одной командой и не затрагивает рабочую базу.

Данные СУБД хранятся в именованном томе \textit{Docker}. Пересборка и перезапуск контейнеров их не затрагивают, а резервные копии в отдельном каталоге защищают от повреждения самого тома.
